% DisplayBoardModule
\section{【中級】DisplayBoardModule — 入出力統合 + ライブラリ内部管理}

\begin{patternbox}{DisplayBoardModule で学べるパターン}
\begin{itemize}
    \item 入出力統合モジュール（updateInput + updateOutput の両方をオーバーライド）
    \item 外部ライブラリ（TFT\_eSPI）をモジュール内部にカプセル化する方法
    \item ModuleTimerによるFPS制御（更新頻度の制限）
    \item Data構造体で入力データ（タッチ座標）と出力データ（表示テキスト）を統合管理
\end{itemize}
\end{patternbox}

\subsection{なぜ1モジュールに統合するのか}

TFTディスプレイボードは「タッチ入力」と「LCD出力」の2機能を持つが、
ハードウェア上は同じSPIバスを共有する1枚の基板であり、同じドライバオブジェクト（TFT\_eSPI）で管理される。
このような場合、BleModuleと同じく\textbf{1つのモジュールで入出力の両方をオーバーライド}する方式が適切である。

\begin{pointbox}[入出力統合の判断基準]
「入力と出力を別モジュールに分けるか、1モジュールに統合するか」は、
\textbf{ハードウェアの単位}で判断する。
同一基板・同一ドライバで管理されるデバイスは統合し、
物理的に独立したデバイスは別モジュールにする。
BleModuleの送受信統合と同じ考え方。
\end{pointbox}

\subsection{ライブラリの内部管理パターン}

\begin{lstlisting}[style=cppstyle, caption={DisplayBoardModule.h — TFT\_eSPIの前方宣言}]
class TFT_eSPI;  // ヘッダへの依存を防ぐ

class DisplayBoardModule : public IModule {
private:
    DisplayBoardConfig _config;
    TFT_eSPI*          _tft;    // 内部で生成・管理
public:
    DisplayBoardModule(const DisplayBoardConfig& config);
    ~DisplayBoardModule();
    bool init()                          override;
    void updateInput(SystemData& data)   override;
    void updateOutput(SystemData& data)  override;
    void deinit()                        override;
};
\end{lstlisting}

\begin{lstlisting}[style=cppstyle, caption={DisplayBoardModule.cpp — コンストラクタでTFT\_eSPIを生成}]
#include <TFT_eSPI.h>

DisplayBoardModule::DisplayBoardModule(const DisplayBoardConfig& config)
    : _config(config), _tft(new TFT_eSPI()) {}

DisplayBoardModule::~DisplayBoardModule() {
    delete _tft;
}
\end{lstlisting}

\begin{pointbox}[外部バスパターンとの違い]
Mpu6500Moduleは\texttt{main.cpp}でI2Cバスを生成し、コンストラクタ引数で渡した（外部バスパターン）。
DisplayBoardModuleではTFT\_eSPIが内部でSPIバスを管理するため、外部から渡す必要がない。
SPIピン設定は\texttt{platformio.ini}のbuild\_flagsでコンパイル時に決定される。
\end{pointbox}

\begin{cautionbox}[既存ライブラリの制約]
TFT\_eSPIやLovyanGFXのように、ライブラリ内部でバスインスタンスを保持するライブラリでは、
アーキテクチャの「外部からバスポインタを渡す」パターンをそのまま適用できない。
この場合、ライブラリごとモジュール内部にカプセル化し、Config構造体や
build\_flagsで設定を渡す方式で対応する。
\end{cautionbox}

\subsection{Data構造体の入出力統合}

\begin{lstlisting}[style=cppstyle, caption={DisplayBoardModule.h — Data}]
struct DisplayBoardData {
    // 出力: 表示テキスト（ロジックフェーズが書き込む）
    char line1[48] = "";
    char line2[48] = "";
    char line3[48] = "";
    char line4[48] = "";
    char line5[48] = "";       // 黄色テキスト
    // 入力: タッチ状態（入力フェーズで更新）
    bool     touchPressed = false;
    uint16_t touchX       = 0; // ピクセル座標
    uint16_t touchY       = 0;
};
\end{lstlisting}

\begin{pointbox}[入力と出力のデータを1つのData構造体に]
\texttt{DisplayBoardData}はタッチ入力（\texttt{touchPressed}等）と
表示出力（\texttt{line1}等）の両方を含む。
入力フェーズでタッチ座標が書き込まれ、ロジックフェーズでそれを参照して
表示テキストを設定し、出力フェーズで画面に描画する。
BleModuleの\texttt{BleData}（受信データ+送信データ）と同じ構成。
\end{pointbox}

\subsection{ロジックフェーズでの表示内容設定}

\begin{lstlisting}[style=cppstyle, caption={main.cpp — applyPattern()での表示設定}]
void applyPattern(SystemData& data) {
    // MPU-6500データを表示用にフォーマット
    if (data.mpu.isValid) {
        snprintf(data.display.line1, sizeof(data.display.line1),
                 "Ax:%.2f Ay:%.2f Az:%.2f",
                 data.mpu.accelX, data.mpu.accelY, data.mpu.accelZ);
    }
    // タッチ座標を表示用にフォーマット
    if (data.display.touchPressed) {
        snprintf(data.display.line5, sizeof(data.display.line5),
                 "Touch: %3d, %3d",
                 data.display.touchX, data.display.touchY);
    }
}
\end{lstlisting}

\begin{pointbox}[3フェーズモデルのデータフロー]
入力フェーズでDisplayBoardModuleが\texttt{data.display.touchPressed/X/Y}に値を書き込み、
ロジックフェーズでMPUデータやタッチ座標を文字列に変換して\texttt{data.display.line1}等に書き込み、
出力フェーズでDisplayBoardModuleが画面に描画する。
入力→ロジック→出力の典型的なデータフローの実例。
\end{pointbox}

% ─── API仕様 ───────────────────────────────
\subsection{API仕様}


\begin{apibox}{DisplayBoardModule --- TFT表示ボードモジュール}
\textbf{定義ファイル:} \texttt{lib/DisplayBoardModule/DisplayBoardModule.h / .cpp}\\
\textbf{分類:} 入出力モジュール（\texttt{updateInput()} + \texttt{updateOutput()}）\\
\textbf{対象デバイス:} ST7789 TFT LCD + XPT2046 タッチパネル

同一基板上のLCD表示（出力）とタッチ読み取り（入力）を統合管理するモジュール。
TFT\_eSPIライブラリを内部で保持し、外部からのドライバ注入は不要。
\end{apibox}

\subsubsection{機能概要}

\begin{itemize}
    \item ST7789 TFT LCDへの5行テキスト表示
    \item XPT2046タッチパネルの座標読み取り
    \item 画面更新周期の制御（\texttt{ModuleTimer}ベース）
    \item バックライトPWM制御
    \item TFT\_eSPIインスタンスの内部管理（new/delete）
    \item \texttt{deinit()}によるリソース解放
\end{itemize}

% --- Config ---
\subsubsection{Config構造体}

\begin{funcblock}{DisplayBoardConfig}
LCD表示のパラメータを定義する。SPI/タッチのピン番号は\texttt{platformio.ini}の\texttt{build\_flags}で設定する。

\begin{longtable}{l l l p{4.5cm}}
\toprule
\textbf{メンバ} & \textbf{型} & \textbf{制約・既定値} & \textbf{説明} \\
\midrule
\endhead
\texttt{rotation}        & \texttt{uint8\_t}  & 0--3（1=横向き） & 画面回転方向 \\
\texttt{blPin}           & \texttt{int8\_t}   & -1で常時ON       & バックライト制御GPIOピン \\
\texttt{updateIntervalMs} & \texttt{uint32\_t} & 例: 100         & 画面更新周期 [ms]。例: 100ms = 10fps \\
\bottomrule
\end{longtable}
\end{funcblock}

\begin{notebox}[SPI/タッチのピン設定]
TFT\_eSPIはSPIバスを内部で管理するため、外部\texttt{SPIClass}は使用しない。
ピン番号（MOSI/MISO/SCK/CS/DC/RST/タッチCS）は\texttt{platformio.ini}の\texttt{build\_flags}で設定する:
\begin{lstlisting}[style=cppstyle, numbers=none, frame=none, backgroundcolor=\color{white}]
build_flags =
    -DTFT_MOSI=11  -DTFT_MISO=13  -DTFT_SCLK=12
    -DTFT_CS=10    -DTFT_DC=14    -DTFT_RST=15
    -DTOUCH_CS=16
\end{lstlisting}
\end{notebox}

% --- Data ---
\subsubsection{Data構造体}

\begin{funcblock}{DisplayBoardData}
LCD表示内容とタッチ入力状態を保持する。

\begin{longtable}{l l l p{4.0cm}}
\toprule
\textbf{メンバ} & \textbf{型} & \textbf{初期値} & \textbf{説明} \\
\midrule
\endfirsthead
\toprule
\textbf{メンバ} & \textbf{型} & \textbf{初期値} & \textbf{説明} \\
\midrule
\endhead
\midrule
\multicolumn{4}{r}{\small 次ページに続く} \\
\endfoot
\bottomrule
\endlastfoot
\multicolumn{4}{l}{\textbf{--- 出力: 表示内容（ロジックフェーズで設定） ---}} \\
\texttt{line1[48]}  & \texttt{char[]}    & \texttt{""} & 1行目テキスト \\
\texttt{line2[48]}  & \texttt{char[]}    & \texttt{""} & 2行目テキスト \\
\texttt{line3[48]}  & \texttt{char[]}    & \texttt{""} & 3行目テキスト \\
\texttt{line4[48]}  & \texttt{char[]}    & \texttt{""} & 4行目テキスト \\
\texttt{line5[48]}  & \texttt{char[]}    & \texttt{""} & 5行目テキスト（タッチ座標表示等） \\
\midrule
\multicolumn{4}{l}{\textbf{--- 入力: タッチ状態（入力フェーズで更新） ---}} \\
\texttt{touchPressed} & \texttt{bool}     & \texttt{false} & タッチ押下状態 \\
\texttt{touchX}       & \texttt{uint16\_t} & \texttt{0}    & タッチX座標 [px] \\
\texttt{touchY}       & \texttt{uint16\_t} & \texttt{0}    & タッチY座標 [px] \\
\end{longtable}
\end{funcblock}

% --- メソッド ---
\subsubsection{メソッド}

\begin{funcblock}{DisplayBoardModule(const DisplayBoardConfig\& config)}
\begin{tabularx}{\textwidth}{l X}
\toprule
\textbf{項目} & \textbf{内容} \\
\midrule
引数   & \texttt{config} --- \texttt{DisplayBoardConfig}へのconst参照 \\
動作   & Configを\texttt{\_config}にコピー。\texttt{TFT\_eSPI}インスタンスを\texttt{new}で生成 \\
\bottomrule
\end{tabularx}
\end{funcblock}

\begin{funcblock}{\textasciitilde DisplayBoardModule()}
\begin{tabularx}{\textwidth}{l X}
\toprule
\textbf{項目} & \textbf{内容} \\
\midrule
動作   & 内部の\texttt{TFT\_eSPI}インスタンスを\texttt{delete}で解放する \\
\bottomrule
\end{tabularx}
\end{funcblock}

\begin{funcblock}{bool init()}
\begin{tabularx}{\textwidth}{l X}
\toprule
\textbf{項目} & \textbf{内容} \\
\midrule
戻り値   & \texttt{bool} --- 常に\texttt{true} \\
動作     & TFT初期化（\texttt{tft.init()}）、画面回転設定、背景色クリア、バックライトON、タッチキャリブレーション \\
\bottomrule
\end{tabularx}
\end{funcblock}

\begin{funcblock}{void updateInput(SystemData\& data)}
\begin{tabularx}{\textwidth}{l X}
\toprule
\textbf{項目} & \textbf{内容} \\
\midrule
書き込み先     & \texttt{data.display}（タッチ関連メンバ） \\
タイミング制御 & 毎ループ実行 \\
動作           & タッチパネルの押下状態を取得し、\texttt{touchPressed}・\texttt{touchX}・\texttt{touchY}を更新 \\
\bottomrule
\end{tabularx}
\end{funcblock}

\begin{funcblock}{void updateOutput(SystemData\& data)}
\begin{tabularx}{\textwidth}{l X}
\toprule
\textbf{項目} & \textbf{内容} \\
\midrule
読み取り元     & \texttt{data.display}（テキスト関連メンバ） \\
タイミング制御 & \texttt{ModuleTimer}で\texttt{updateIntervalMs}ごとに実行 \\
動作           & 画面をクリアし、\texttt{line1}--\texttt{line5}のテキストを描画する \\
\bottomrule
\end{tabularx}
\end{funcblock}

\begin{funcblock}{void deinit()}
\begin{tabularx}{\textwidth}{l X}
\toprule
\textbf{項目} & \textbf{内容} \\
\midrule
動作   & バックライトOFF、画面クリア \\
\bottomrule
\end{tabularx}
\end{funcblock}

% --- 使用例 ---
\subsubsection{使用例}

\begin{lstlisting}[style=cppstyle, caption={DisplayBoardModuleの使用例}]
// ProjectConfig.h
const DisplayBoardConfig DISPLAY_BOARD_CONFIG = {
    .rotation        = 1,    // 横向き
    .blPin           = 46,   // バックライトピン
    .updateIntervalMs = 100, // 10fps
};

// ロジックフェーズでの使用
void applyPattern(SystemData& data) {
    // センサーデータを表示
    snprintf(data.display.line1, 48,
             "Accel: %.2f g", data.mpu.accelX);
    snprintf(data.display.line2, 48,
             "Battery: %.1f V", data.battery.voltage);
    // タッチ座標の表示
    if (data.display.touchPressed) {
        snprintf(data.display.line5, 48,
                 "Touch: (%d, %d)",
                 data.display.touchX,
                 data.display.touchY);
    }
}
\end{lstlisting}
